(1) 取右为平动正向、逆时针为角加速度正向，地面摩擦力 \(f\) 以向右为正。水平动量式 \(ma=F\cos\alpha+f\)，绕质心转矩式 \(I\dot\omega=rF+Rf\)；接触点不滑动给出 \(a=-R\dot\omega\)。这些符号须按题设的出线侧选取；若线绕相反侧，驱动转矩符号也反向。
(2) 将 \(f=ma-F\cos\alpha\) 代入转矩式，有 \(-Ia/R=rF+Rma-RF\cos\alpha\)，于是 \(a=FR(R\cos\alpha-r)/(I+mR^2)\)。当 \(\cos\alpha>r/R\) 时向右加速，\(\cos\alpha<r/R\) 时向左加速，临界角 \(\alpha_c=\arccos(r/R)\)；方向反转并不表示拉力反向。
(3) 由上式 \(f=F[mR(R\cos\alpha-r)/(I+mR^2)-\cos\alpha]\)，静摩擦力的方向可正可负，不能预设恒向左。垂直平衡给 \(N=mg-F\sin\alpha\)；必须有 \(N>0\) 才能保持接触。若接触成立，纯滚动必要条件是 \(\mu_s\ge |f|/N\)，等号表示静摩擦已达到库仑极限。
(4) 接触点瞬时静止，所以地面静摩擦的功率为零。拉力作用点水平速度为 \(v+r\omega\) 对应抽线方向投影，合外力功率可写 \(P=F(\cos\alpha\,v+r\omega)\)；因为 \(\omega=-v/R\)，\(P=Fv(\cos\alpha-r/R)\)。总动能 \(K=(m+I/R^2)v^2/2\)，其导数 \((m+I/R^2)va\) 与代入上问 \(a\) 后的功率一致。
